% =========================
% 4_overview.tex
% =========================
\section{System Overview}
\label{sec:overview}

HALO is an IMU-to-text semantic alignment model trained across heterogeneous HAR datasets for open-set, zero-shot activity recognition, both offline and as a live stream. The system addresses the design requirements identified in Section~\ref{sec:motivation} through a heterogeneity-aware encoder trained in a single language-alignment stage. Figure~\ref{fig:overview} illustrates the end-to-end pipeline.


\begin{figure}[t]
    \centering
    \setlength{\abovecaptionskip}{2pt}
    \setlength{\belowcaptionskip}{-4pt}
  \includegraphics[width=\columnwidth]{figures/fig_overview.png}
  \caption{System overview of HALO. The sensing heterogeneity is addressed through heterogeneity-aware pretraining of the IMU encoder (Stage~1), followed by synonym-aware soft contrastive learning to align sensor features with text semantics (Stage~2), enabling generalizable and open-set activity recognition. Zero-shot inference is achieved by swapping label banks at runtime without retraining.  }
  \label{fig:overview}
\end{figure}

% System overview of HALO. Sensing heterogeneity is addressed through heterogeneity-aware pretraining of the IMU encoder (Stage 1), followed by synonym-aware contrastive learning to align sensor features with text semantics (Stage 2), enabling generalizable and open-set activity recognition. Zero-shot inference is achieved by swapping label banks at runtime without retraining.

\textbf{Heterogeneity-aware encoder.} The IMU encoder maps variable-rate, variable-channel sessions to per-patch embeddings through a physical-frequency filterbank tokenizer (rate-invariant by construction), channel-independent processing, and language-based sensor conditioning; a dual-branch transformer with rotary position embeddings over physical time then fuses them (Section~\ref{sec:heterogeneity}).

\textbf{Single-stage language alignment.} The encoder is aligned---from scratch, in a single stage---with a frozen Sentence-BERT text encoder via synonym-aware soft contrastive learning that derives soft target distributions from text--text similarity, resolving label conflicts across datasets and mapping both modalities into a shared 384-dimensional space (Section~\ref{sec:alignment_design}).

\textbf{Offline and streaming from one model.} Because the tokenizer is per-patch (no cross-patch statistics) and positions are relative (rotary over physical time), and because temporal attention is trained under both full and causal masks, the same weights run two ways: \emph{offline} over a finished recording (full context, pooled to a session label) and \emph{online} as a causal, KV-cached stream that emits a prediction per patch (Section~\ref{sec:streaming}).

\textbf{Inference.} At test time, recognition is cosine-similarity retrieval: given a patch (or pooled) embedding, we select the label embedding with maximum similarity, supporting arbitrary candidate inventories including open-vocabulary settings. No per-dataset classifier is needed. \STALE{The v1 model contains ${\sim}$35M trainable parameters (8-layer dual-branch IMU encoder ${\sim}$19M, alignment and projection heads ${\sim}$16M)} plus a single frozen Sentence-BERT backbone (all-MiniLM-L6-v2, 22.7M) shared for both channel conditioning and text alignment. \OUTDATED{The v2 parameter breakdown changes: the physical-Hz filterbank tokenizer replaces the CNN/spectral extractor, and the learnable label bank is dropped in favor of mean pooling.} At deployment, the sensor description string is specified once when the app is configured; the corresponding text embeddings are computed once and cached, adding zero per-inference cost. Likewise, all candidate label embeddings can be precomputed offline.

% Figure~\ref{fig:overview} illustrates the end-to-end pipeline; Figures~\ref{fig:encoder} and~\ref{fig:alignment} detail the encoder and alignment components respectively.
